\documentclass[12pt]{ximera}
\title{Practice Exam 3}
\begin{document}
\begin{abstract}
Practice for Exam 3: properties of discrete and continuous division, which methods can
leave every player better than fair, the method of markers, a three-round lone-chooser
division of a sandwich, and sealed bids across four choir sections.
\end{abstract}
\maketitle

\section*{Practice Exam 3}

\begin{problem}
Assess whether each of the following claims is true or false.

\begin{prompt}
\textbf{(a)} Claim: in a system with discrete assets, it is always possible to achieve a
perfectly equal division.

\begin{multipleChoice}
\choice{True}
\choice[correct]{False}
\end{multipleChoice}

\begin{feedback}
False. Discrete assets cannot be subdivided, so there may simply be no way to split
them evenly.

The simplest counterexample: two players and one car. Whoever does not get the car gets
nothing, and no rearrangement fixes that. This is exactly why the discrete methods bring
in \emph{money} --- the compensation in sealed bids exists to patch the gaps that
indivisible objects leave behind.
\end{feedback}
\end{prompt}

\begin{prompt}
\textbf{(b)} Claim: in a system with continuous assets and $N$ total players, it is
always possible to achieve a fair division by splitting every asset into $N$ parts, one
of each going to every player.

\begin{multipleChoice}
\choice[correct]{True}
\choice{False}
\end{multipleChoice}

\begin{feedback}
True. If every asset is divided into $N$ equal parts and each player receives one part
of each asset, then each player receives exactly $\frac1N$ of every asset --- and
therefore exactly $\frac1N$ of the total value \emph{in their own value system},
whatever that system happens to be.

This works regardless of how wildly the players' preferences differ, which is what makes
it a guaranteed fair division. It is not, however, a very good one: nobody gets more
than the bare minimum, and the methods we study exist precisely to do better than this
by exploiting the players' differing tastes.
\end{feedback}
\end{prompt}
\end{problem}

\begin{problem}
Determine all applicable answers to the following.

\begin{prompt}
\textbf{(a)} In which of the following fair division methods is it possible for
\emph{every} player to come out with \emph{greater} than their fair share?

\begin{selectAll}
\choice{Lone-Divider Method}
\choice[correct]{Sealed Bids Method}
\choice{Lone-Chooser Method}
\choice[correct]{Method of Markers}
\end{selectAll}

\begin{feedback}
\textbf{Sealed Bids} and the \textbf{Method of Markers}.

In Sealed Bids, everyone is brought to exactly a fair share and then the surplus is
split equally, leaving every player strictly ahead. In the Method of Markers, the
leftover items are distributed after everyone already holds a fair share.

The divider-based methods cannot manage it, because the divider always ends up with
exactly $\frac1N$ --- see part (b).
\end{feedback}
\end{prompt}

\begin{prompt}
\textbf{(b)} In which of the following fair division methods does at least one player
\emph{always} come out with exactly their \emph{minimum} fair share?

\begin{selectAll}
\choice[correct]{Lone-Divider Method}
\choice{Sealed Bids Method}
\choice[correct]{Lone-Chooser Method}
\choice{Method of Markers}
\end{selectAll}

\begin{hint}
Think about the player who does the cutting in each method. What is every piece worth to
them?
\end{hint}

\begin{feedback}
\textbf{Lone-Divider} and \textbf{Lone-Chooser} --- in both, the first divider is stuck
at exactly a fair share.

In \emph{Lone-Divider}, the divider cuts the asset into $N$ pieces they consider
identical, so whichever piece they receive is worth exactly $\frac1N$ to them.

In \emph{Lone-Chooser}, the same logic applies one step at a time. With three players,
player 1 cuts the asset into two halves worth $\frac12$ each to them, and keeps one.
They then cut that half into three pieces worth $\frac16$ each to them, and keep two.
That is $\frac26 = \frac13$ --- exactly a fair share, no matter what the other players
do.

Homework 8 illustrates this: group B, who cut first, ended up with exactly
$33.\overline{3}\%$ while groups A and C came out ahead.
\end{feedback}
\end{prompt}
\end{problem}

\begin{problem}
In the following discrete system with three players, we use the method of markers. Each
player has placed two markers in the array of $13$ items.

\begin{center}
\includegraphics[width=\textwidth]{markersArray.png}
\end{center}

The markers fall as follows:
\[
\text{Player A}: 4,\ 9 \qquad \text{Player B}: 3,\ 8 \qquad \text{Player C}: 4,\ 7
\]

\begin{prompt}
Which items will go to each player, and which items will be surplus?

Player $\answer{B}$ gets items 1--$\answer{3}$.

\smallskip
Player $\answer{C}$ gets items 4--$\answer{7}$.

\smallskip
Player $\answer{A}$ gets items $\answer{10}$--13.

\smallskip
The surplus is items $\answer{8}$ and $\answer{9}$.

\begin{feedback}
\textbf{Round 1.} First markers are $A_1=4$, $B_1=3$, $C_1=4$. B's is leftmost, so
\textbf{B takes items 1--3}.

\textbf{Round 2.} Among A and C, the second markers are $A_2=9$ and $C_2=7$. C's is
leftmost, so \textbf{C takes items 4--7}.

\textbf{A} is the last player left and takes everything past their final marker at $9$:
\textbf{items 10--13}.

\textbf{Surplus: items 8 and 9}, which fall between the end of C's segment and the start
of A's.

Note that every player does at least as well as they asked for: B's first segment was
exactly $1$--$3$, C's second segment ran $5$--$7$ and C received $4$--$7$, and A's final
segment began at $10$.
\end{feedback}
\end{prompt}
\end{problem}

\begin{problem}
Jake, Adam, and Joel want to share an $18$-inch half-veggie / half-roast-beef sandwich,
using the \emph{lone-chooser} method, each time making only the minimum necessary
vertical cuts. After a mini-tournament of rock-paper-scissors, Adam will be the first player, Jake the
second, and Joel the third.
\begin{itemize}
\item Adam likes roast beef three times as much as veggie;
\item Jake is strictly vegetarian, so doesn't value roast beef at all;
\item Joel likes veggie and roast beef equally.
\end{itemize}

\begin{prompt}
\textbf{(a)} Complete the chart with each player's value system, as fractions.

Adam --- Veggie: $\answer{1/4}$, \quad Beef: $\answer{3/4}$

\smallskip
Jake --- Veggie: $\answer{1}$, \quad Beef: $\answer{0}$

\smallskip
Joel --- Veggie: $\answer{1/2}$, \quad Beef: $\answer{1/2}$

\begin{feedback}
Adam's $1:3$ ratio gives $\frac14$ and $\frac34$. Jake places all of his value on the
veggie half: $1$ and $0$. Joel splits evenly: $\frac12$ and $\frac12$.
\end{feedback}
\end{prompt}

\begin{prompt}
\textbf{(b, Round 1)} Adam divides the sandwich into two shares of equal value to him,
and Jake selects one.

Adam's two shares: all of V $+\ \answer{1/3}$ of B, \quad and \quad $\answer{2/3}$ of B.

\smallskip
Jake selects the share containing the veggie, which is worth $\answer{1}$ of the whole
to him.

\begin{hint}
All of the veggie half is worth only $\frac14$ to Adam --- less than the $\frac12$ he
needs --- so his cut has to reach into the beef.
\end{hint}

\begin{feedback}
Adam needs two shares worth $\frac12$ each. All of V is worth only $\frac14$, so his
share must also include a fraction $x$ of the beef:
\[ \tfrac14 + x\cdot\tfrac34 = \tfrac12 \quad\Longrightarrow\quad x = \tfrac13. \]
So Adam cuts at all of V $+\frac13$ of B, leaving $\frac23$ of B. (On an $18$-inch
sandwich, that is a cut $3$ inches into the beef.)

Jake values only the veggie, so the first share is worth the \emph{entire} value of the
sandwich to him --- a full $1$, or $100\%$ --- while the second is worth $0$. He takes
the first.
\end{feedback}
\end{prompt}

\begin{prompt}
\textbf{(b, Round 2)} Each of Adam and Jake now cuts their share into three pieces of
equal value to themselves.

Adam's three pieces are each $\answer{2/9}$ of B.

\smallskip
Jake's three pieces: $\frac13$ of V, \quad $\frac13$ of V, \quad and
$\answer{1/3}$ of V $+\ \frac13$ of B.

\smallskip
Joel then selects one piece from each. His best total is worth $\answer{4/9}$ of the
whole to him.

\begin{hint}
Adam's share is all beef, so equal value means equal size. Jake values only veggie, so
his cuts must divide the \emph{veggie} into thirds --- the beef he is holding is worth
nothing to him and simply rides along with the last piece.
\end{hint}

\begin{feedback}
\textbf{Adam} holds $\frac23$ of B. Since it is all beef, three equal-value pieces are
three equal-size pieces: $\frac29$ of B each.

\textbf{Jake} holds all of V $+\frac13$ of B, which he values entirely by the veggie. Two
cuts through the veggie at the one-third and two-thirds marks give him three pieces he
values equally, and the leftover beef rides along on the last piece:
\[ \tfrac13 \text{ of V}, \qquad \tfrac13 \text{ of V}, \qquad \tfrac13 \text{ of V} + \tfrac13 \text{ of B}. \]

\textbf{Joel} ($V=\frac12$, $B=\frac12$) evaluates:
\begin{itemize}
\item each of Adam's pieces: $\frac29 \cdot \frac12 = \frac19$ --- all three are equal;
\item Jake's first two pieces: $\frac13\cdot\frac12 = \frac16$ each;
\item Jake's third piece: $\frac16 + \frac13\cdot\frac12 = \frac16+\frac16 = \frac13$.
\end{itemize}
He takes any of Adam's pieces plus Jake's third piece, for
\[ \tfrac19 + \tfrac13 = \tfrac19+\tfrac39 = \tfrac49 \approx 44.4\%. \]

\textbf{Final shares}, and what each is worth to its owner:
\begin{itemize}
\item \textbf{Adam}: $\frac49$ of B, worth $\frac49\cdot\frac34 = \frac13$ --- exactly a
      fair share, as the first divider always gets.
\item \textbf{Jake}: $\frac23$ of V, worth $\frac23 \approx 66.7\%$.
\item \textbf{Joel}: $\frac13$ of V $+\frac59$ of B, worth $\frac49 \approx 44.4\%$.
\end{itemize}
All the veggie and all the beef are accounted for, and everyone clears $\frac13$.
\checkmark
\end{feedback}
\end{prompt}
\end{problem}

\begin{problem}
A community chamber choir is hosting a gala and need to distribute
responsibilities across Sopranos, Altos, Tenors, and Basses. Each section leader has
submitted sealed bids indicating how much they mind being assigned each task.

\begin{center}
\begin{tabular}{|r|r|r|r|r|}\hline
 & \textbf{Sopranos} & \textbf{Altos} & \textbf{Tenors} & \textbf{Basses}\\\hline
Event planning & $-\$80$ & $-\$50$ & $-\$100$ & $-\$70$ \\
Advertisement & $-\$20$ & $-\$30$ & $-\$50$ & $-\$40$ \\
Riser moving & $-\$110$ & $-\$100$ & $-\$40$ & $-\$60$ \\
Baking desserts & $-\$10$ & $-\$30$ & $-\$90$ & $-\$50$ \\
Clean-up crew & $-\$20$ & $-\$10$ & $-\$20$ & $-\$20$ \\\hline
\end{tabular}
\end{center}

\begin{prompt}
\textbf{(a)} Which section is assigned each task?

Event planning: $\answer{Altos}$ \quad Advertisement: $\answer{Sopranos}$ \quad
Riser moving: $\answer{Tenors}$

\smallskip
Baking desserts: $\answer{Sopranos}$ \quad Clean-up crew: $\answer{Altos}$

\begin{feedback}
Each task goes to the highest bidder --- here, the least negative bid:
\[
\begin{array}{ll}
\text{Event planning} & \text{Altos } (-50)\\
\text{Advertisement} & \text{Sopranos } (-20)\\
\text{Riser moving} & \text{Tenors } (-40)\\
\text{Baking desserts} & \text{Sopranos } (-10)\\
\text{Clean-up crew} & \text{Altos } (-10)
\end{array}
\]
Note that the Basses win nothing --- they are never the least reluctant on any task.
\end{feedback}
\end{prompt}

\begin{prompt}
\textbf{(b)} Compute each section's total value and fair share.

Total value --- Sopranos: $\answer{-240}$ \quad Altos: $\answer{-220}$ \quad
Tenors: $\answer{-300}$ \quad Basses: $\answer{-240}$

\smallskip
Fair share --- Sopranos: $\answer{-60}$ \quad Altos: $\answer{-55}$ \quad
Tenors: $\answer{-75}$ \quad Basses: $\answer{-60}$

\begin{feedback}
Summing each column and dividing by $4$:
\[
\begin{array}{lll}
\text{Sopranos}: & -240 & \to -\$60\\
\text{Altos}: & -220 & \to -\$55\\
\text{Tenors}: & -300 & \to -\$75\\
\text{Basses}: & -240 & \to -\$60
\end{array}
\]
\end{feedback}
\end{prompt}

\begin{prompt}
\textbf{(c)} Compute each section's allotment value and initial settlement. Enter the
settlement as a positive number for money paid in, negative for money received.

Allotment --- Sopranos: $\answer{-30}$ \quad Altos: $\answer{-60}$ \quad
Tenors: $\answer{-40}$ \quad Basses: $\answer{0}$

\smallskip
Initial settlement --- Sopranos: $\answer{30}$ \quad Altos: $\answer{-5}$ \quad
Tenors: $\answer{35}$ \quad Basses: $\answer{60}$

\begin{feedback}
Allotment values: Sopranos have Advertisement and Baking, $-20-10 = -\$30$; Altos have
Event planning and Clean-up, $-50-10 = -\$60$; Tenors have Riser moving, $-\$40$;
Basses have nothing, $\$0$.

Subtracting each fair share:
\[
\begin{array}{llll}
\text{Sopranos}: & -30 - (-60) = +30 & \text{pays } \$30\\
\text{Altos}: & -60 - (-55) = -5 & \text{receives } \$5\\
\text{Tenors}: & -40 - (-75) = +35 & \text{pays } \$35\\
\text{Basses}: & 0 - (-60) = +60 & \text{pays } \$60
\end{array}
\]
The Altos are the only section to receive money at this stage --- their two tasks add up
to more unpleasantness than their fair share.
\end{feedback}
\end{prompt}

\begin{prompt}
\textbf{(d)} What is the surplus, and what is each section's share of it?

Surplus: $\$\answer{120}$ \qquad Each section's share: $\$\answer{30}$

\begin{feedback}
Money in is $30+35+60 = \$125$; money out is $\$5$. The surplus is
\[ \$125 - \$5 = \$120, \]
which splits four ways as $\$30$ each.
\end{feedback}
\end{prompt}

\begin{prompt}
\textbf{(e)} Give each section's final settlement. Enter a positive number for money
paid, negative for money received.

Sopranos: $\answer{0}$ \quad Altos: $\answer{-35}$ \quad Tenors: $\answer{5}$ \quad
Basses: $\answer{30}$

\begin{feedback}
Subtract each section's \$30 share of the surplus from their initial settlement:
\[
\begin{array}{ll}
\text{Sopranos}: & 30 - 30 = \$0 \text{ --- pays nothing, does Advertisement and Baking}\\
\text{Altos}: & -5 - 30 = -\$35 \text{ --- receives \$35, does Event planning and Clean-up}\\
\text{Tenors}: & 35 - 30 = \$5 \text{ --- pays \$5, does Riser moving}\\
\text{Basses}: & 60 - 30 = \$30 \text{ --- pays \$30, does no tasks}
\end{array}
\]
The cash balances: $0 - 35 + 5 + 30 = 0$. \checkmark

Every section finishes \$30 better off than a fair share in its own valuation. The
Basses do no work at all but pay \$30 for the privilege --- which they consider a
bargain, since the tasks they avoided were worth $-\$240$ to them.
\end{feedback}
\end{prompt}
\end{problem}

\end{document}
